% ============================================================
% 00-abstract.tex — V13 Rewritten Abstract
% ============================================================
\begin{abstract}
University counseling spans relational curricula, cohort-specific fees, exemption policies, and narrative regulations, each requiring distinct evidence representations and computational treatments. Whether decomposing this heterogeneous task into multiple specialized agents improves end-to-end reliability over a single agent with equivalent capabilities has not been directly evaluated in prior work. We present CTU-Chat, a supervisor-routed multi-agent system that assigns queries to four bounded specialists with scoped tool privileges and representation-specific evidence paths: Neo4j graph traversal for curricula, structured lookup with deterministic calculation for financial data, and hybrid document retrieval with neural reranking for regulations. We evaluate CTU-Chat through three experimental scenarios on a 100-query held-out benchmark: (1)~an architecture-level comparison against a functionally matched single-agent baseline and a routed generic-worker variant under identical model, knowledge, and tool capabilities; (2)~architectural ablations that individually remove specialist policies, tool isolation, heterogeneous evidence allocation, deterministic calculation, and route repair; and (3)~a semantic-neighbor tool-ambiguity stress test across registries of 11, 31, and 51 tools. %
% [PLACEHOLDER: Insert primary quantitative results after experiments]
These findings suggest that domain-scoped specialization and heterogeneous evidence allocation can support institutional counseling reliability under the evaluated conditions, with measurable trade-offs in coordination latency and token cost.

\keywords{Multi-Agent Systems \and Supervisor Routing \and Architectural Ablation \and Heterogeneous Knowledge Allocation \and Retrieval-Augmented Generation \and University Counseling}
\end{abstract}
